\chapter{Separable equations and their disguises}

\begin{orient}
\textbf{What this chapter is for.} Separability is the cheapest thing that can happen to a first-order equation, and it happens far more often than it appears to, because problems are built to hide it. The visible form $y'=g(x)h(y)$ is rare; the common situation is an equation that becomes separable after a grouping, an identity, or one of a very short list of substitutions. That list is the content of this chapter: $y=vx$ when the right side depends only on $y/x$, $u=ax+by+c$ when it depends only on that combination, and a handful of trigonometric substitutions that rationalise a stubborn side. Alongside the substitutions runs a discipline that matters more than any of them. Dividing by $h(y)$ destroys the solutions with $h(y)=0$, and those constant solutions are frequently the equilibria that the problem is actually about. Recording them is part of the method, not an afterthought, and a solution set that omits them is wrong rather than incomplete. By the end you should be able to separate reliably, name the lost solutions before you lose them, and recognise the four or five disguises that account for nearly every hidden separable equation you will meet.
\end{orient}

\section{The protocol}

Separation is four steps, and skipping any of them is where the marks go. Write the equation as $\dv{y}{x}=g(x)h(y)$. Find and record every root of $h$, since each root $y_0$ gives a constant solution $y\equiv y_0$. Divide by $h(y)$ and integrate both sides, putting a single constant on the $x$ side. Then apply the initial condition and, separately, check whether the constant solutions you recorded are recovered by the general formula for some value of the constant, or whether they must be reported separately.

The third step deserves a comment about form. Chasing a constant through an exponential and a logarithm is where most arithmetic errors enter, and the definite form of the integration avoids all of it. Writing $\displaystyle\int_{y_0}^{y}\frac{\dd s}{h(s)}=\int_{x_0}^{x}g(t)\dd t$ builds the initial condition into the statement, so there is no constant to determine and no exponentiation of an unknown $C$ to mishandle. It also makes the domain question visible: the identity holds only as long as the left-hand integral is finite and $h$ does not vanish between $y_0$ and $y$, which is exactly the condition under which the solution exists and is unique.

The fourth step, the constant-solution check, is the one that gets skipped, and it is the one that changes answers. It is worth being explicit about why the roots of $h$ matter. If $h(y_0)=0$ then $y\equiv y_0$ satisfies $y'=g(x)h(y)$ identically, so it is a solution; and dividing by $h(y)$ in step three is a division by zero on that solution, which is why the general formula usually cannot see it. Sometimes the formula recovers it anyway at a limiting value of the constant, and sometimes it does not. Deciding which happened takes one line and is part of the answer.

\tech{The separable machine with an initial condition}{core}
\cue $\dv{y}{x}=g(x)h(y)$ together with data $y(x_0)=y_0$, in any of the equivalent arrangements $M(x)\dd x+N(y)\dd y=0$ or $N(y)y'=M(x)$.
\method Prefer the definite form, which applies the initial condition automatically:
\[\int_{y_0}^{y}\frac{\dd s}{h(s)}=\int_{x_0}^{x}g(t)\dd t.\]
No constant of integration appears and none has to be chased. Solve for $y$ explicitly only if the question needs it; an implicit relation is a complete answer. Before dividing, record every root of $h$: those are constant solutions, and the definite form is valid only on an interval where $h(y)\neq0$.

\begin{worked}
$\dv{y}{x}=\dfrac{4+9y^{2}}{\eu^{2x+1}}$ with $y(-\tfrac12)=0$.

There are no roots of $h(y)=4+9y^{2}$, so nothing is lost. In definite form,
\[\int_{0}^{y}\frac{\dd s}{4+9s^{2}}=\int_{-1/2}^{x}\eu^{-(2t+1)}\dd t.\]
The left side is $\tfrac16\arctan\tfrac{3y}{2}$ and the right is $\tfrac12\bigl(1-\eu^{-(2x+1)}\bigr)$, so
\[y=\frac23\tan\!\Bigl(3\bigl(1-\eu^{-(2x+1)}\bigr)\Bigr).\]
Now read the formula for what it says about the domain. The tangent's argument increases from $0$ toward $3$ as $x$ increases, and it crosses $\tfrac{\pi}{2}$ when $1-\eu^{-(2x+1)}=\tfrac{\pi}{6}$, that is at
\[x^{*}=\tfrac12\Bigl(-1-\ln\bigl(1-\tfrac{\pi}{6}\bigr)\Bigr)=-0.1292526.\]
The solution blows up there and does not exist beyond it, which the differential equation itself never announced. Numerically, $y(-0.2)=3.0205002$ from both the formula and a Runge-Kutta integration, and $y(-0.15)=10.9963$; the agreement holds right up to the asymptote.
\end{worked}

\begin{trap}
Writing $\ln\abs y=\cdots+C$ and exponentiating to $y=\eu^{C}\eu^{\cdots}$ while quietly forgetting three things: the absolute value, which means the new constant $A=\pm\eu^{C}$ ranges over all nonzero reals rather than only positive ones; the value $A=0$, which is not of the form $\pm\eu^{C}$ and which recovers the equilibrium solution $y\equiv0$; and the fact that a solution can cease to exist at finite $x$ even when $g$ and $h$ are perfectly smooth everywhere.
\end{trap}

\begin{seealsob}The separability test; hidden separability by factoring; autonomous recognition and the phase line.\end{seealsob}

The next example is the one to keep in mind whenever the constant-solution question feels pedantic, because there the constants are the whole story: they are the equilibria, they bound every other solution, and the closed form is meaningless without them.

\begin{worked}[Worked example: the lost solutions are the answer]
$\dv{y}{x}=x\bigl(y^{2}-4\bigr)$ with $y(0)=0$.

Roots of $h(y)=y^{2}-4$: $y=2$ and $y=-2$, so $y\equiv2$ and $y\equiv-2$ are constant solutions. Record them, then divide. Partial fractions give $\dfrac{1}{y^{2}-4}=\dfrac{1}{4}\left(\dfrac{1}{y-2}-\dfrac{1}{y+2}\right)$, so
\[\tfrac14\ln\left|\frac{y-2}{y+2}\right|=\frac{x^{2}}{2}+c.\]
At $x=0$, $y=0$, the left side is $\tfrac14\ln1=0$, so $c=0$ and $\dfrac{y-2}{y+2}=-\eu^{2x^{2}}$, the sign fixed by the initial point. Solving,
\[y=\frac{2\bigl(1-\eu^{2x^{2}}\bigr)}{1+\eu^{2x^{2}}}=-2\tanh\bigl(x^{2}\bigr).\]
Check: $y(1.5)=-2\tanh(2.25)=-1.9560522$, which Runge-Kutta reproduces to ten figures. Now the recorded constants pay. The solution decreases from $0$ and approaches $-2$ but never reaches it, because $y\equiv-2$ is itself a solution and solutions of a Lipschitz equation cannot cross. That single sentence gives the long-run behaviour, the monotonicity, and the bound, and it is unavailable to anyone who divided by $y^{2}-4$ without looking. Note also that neither $y\equiv2$ nor $y\equiv-2$ is recovered by the formula for any finite constant; both must be reported separately.
\end{worked}

\begin{keynote}[The constant-solution check is part of the method]
Every division by $h(y)$ is a case split, and the discarded case is a solution. Write the roots down \emph{before} dividing, then check each against the final formula. Three outcomes are possible. The formula recovers the root at some finite constant, in which case nothing was lost. The formula recovers it only in a limit, as with $A\to0$ in $y=A\eu^{kx}$, in which case the root must be added to the solution set by hand. Or the root is a genuinely separate branch through which other solutions pass, in which case uniqueness has failed and the initial-value problem has more than one solution. The standard example of the third is $y'=y^{2/3}$ with $y(0)=0$, which has the solutions $y\equiv0$ and $y=\tfrac{1}{27}x^{3}$ and infinitely many hybrids that sit on zero for a while and then leave. The Lipschitz condition fails exactly at $y=0$, and that is where the lost constant solution was.
\end{keynote}

\section{Separability that has been hidden}

The commonest disguise is not a substitution at all. It is arithmetic: an equation that separates after a grouping, a factorisation, or an elementary identity, presented in expanded form so that no factorisation is visible. This is why the audit begins with ``simplify'' and why the first move on any first-order equation is to factor every coefficient completely.

\tech{Hidden separability by factoring or by an identity}{easy}
\cue The equation looks like a sum, or a bloated rational expression, or carries a radical whose radicand is suspiciously structured. Trigonometric coefficients that are longer than one term are almost always an identity in disguise.
\method Before choosing any method, factor every coefficient completely and simplify with elementary identities. Group four-term sums in pairs. Recognise the standard collapses: $\sin^{2}x+\cos^{2}x+\tan^{2}x=\sec^{2}x$; $1+\tan^{2}x=\sec^{2}x$; $x+2\sqrt{x-1}=\bigl(\sqrt{x-1}+1\bigr)^{2}$; $\eu^{x+y}=\eu^{x}\eu^{y}$, which turns an apparent sum in the exponent into a product and separates immediately.

\begin{worked}
$\dv{y}{x}=\sqrt{x+2\sqrt{x-1}}$ with $y(1)=3$; find $y(5)$.

The radicand is a perfect square: $x+2\sqrt{x-1}=(x-1)+2\sqrt{x-1}+1=\bigl(\sqrt{x-1}+1\bigr)^{2}$. On $x>1$ the inner expression is positive, so the outer square root removes cleanly and
\[\dv{y}{x}=\sqrt{x-1}+1,\qquad y=\tfrac23(x-1)^{3/2}+x+C.\]
From $y(1)=3$ we get $C=2$, so $y(5)=\tfrac23(8)+5+2=\tfrac{37}{3}=12.3333$, confirmed numerically.

A second instance, the four-term sum: $\dv{y}{x}=xy+x+y+1$. Grouped, this is $(x+1)(y+1)$, so $\ln\abs{1+y}=\tfrac{x^{2}}{2}+x+c$ and $y=A\eu^{x^{2}/2+x}-1$. The constant solution $y\equiv-1$ is the case $A=0$. With $y(0)=0$, $A=1$ and $y(1)=\eu^{3/2}-1=3.4816891$.
\end{worked}

A third instance, where the disguise is a rational expression rather than a radical: $\dv{y}{x}=\dfrac{x+xy^{2}}{y+x^{2}y}$ with $y(0)=1$. Both numerator and denominator factor, and the factorisation is the whole problem:
\[\dv{y}{x}=\frac{x\bigl(1+y^{2}\bigr)}{y\bigl(1+x^{2}\bigr)},\qquad \frac{y\dd y}{1+y^{2}}=\frac{x\dd x}{1+x^{2}},\]
so $\ln\bigl(1+y^{2}\bigr)=\ln\bigl(1+x^{2}\bigr)+c$ and $1+y^{2}=C\bigl(1+x^{2}\bigr)$. From $y(0)=1$, $C=2$ and $y^{2}=1+2x^{2}$, so $y(2)=3$, confirmed numerically. Expanded, the original right-hand side is a ratio of two-term sums and looks like nothing; factored, it is a product of a pure-$x$ and a pure-$y$ piece. The habit worth building is to factor numerator and denominator \emph{separately} before judging separability, since a common factor between them would also have simplified the problem.

\begin{trap}
Dropping the absolute value when un-squaring. Strictly $\sqrt{(\sqrt{x-1}+1)^{2}}=\abs{\sqrt{x-1}+1}$, which equals $\sqrt{x-1}+1$ only where the inner expression is nonnegative; here that is all of $x\geq1$, but the domain must be stated rather than assumed. The general form of the trap is that $\sqrt{u^{2}}=\abs u$, and problems are built around radicands that are perfect squares of expressions which change sign.
\end{trap}

\begin{seealsob}The separability test; the classification audit; simplification before any method is chosen.\end{seealsob}

A short catalogue of the collapses worth recognising on sight, since they account for most hidden separability in practice. An exponential of a sum, $\eu^{x+y}=\eu^{x}\eu^{y}$, separates immediately and is the commonest disguise of all: $y'=\eu^{x+y}$ looks like a sum and is a product. A four-term sum with a shared binomial, $xy+x+y+1=(x+1)(y+1)$, is the second commonest. A radicand of the form $u\pm2\sqrt{u-1}$ is $\bigl(\sqrt{u-1}\pm1\bigr)^{2}$. A trigonometric coefficient of more than one term is nearly always a Pythagorean identity: $\sin^{2}x+\cos^{2}x+\tan^{2}x=\sec^{2}x$, and $\dfrac{1-\cos2x}{2}=\sin^{2}x$. And a logarithm of a product or a power should be expanded before anything else, since $\ln(x^{2}y)=2\ln x+\ln y$ turns a mixed expression into a sum whose separability is then visible.

The first genuine substitution handles the case where the right side is not a product but depends on $x$ and $y$ only through their ratio. This is the degree-zero homogeneous family, and the substitution that separates it is universal within that family: it never fails, and it always produces the same shape of separated equation.

\tech[Substitution y equals v x for homogeneous equations]{Substitution $y=vx$ for degree-zero homogeneous equations}{medium}
\cue The right side passes the degree-zero test $f(tx,ty)=f(x,y)$, or equivalently every term of $M$ and $N$ in $M\dd x+N\dd y=0$ has the same total degree. Writing $f$ in terms of $y/x$ and finding that $x$ disappears is the confirmation.
\method Put $y=vx$, so that $y'=v+x\dv{v}{x}$. Every degree-zero equation $y'=F(v)$ becomes
\[x\dv{v}{x}=F(v)-v,\qquad\text{that is}\qquad \frac{\dd v}{F(v)-v}=\frac{\dd x}{x},\]
which is separable in every case. The roots of $F(v)-v=0$ are lost solutions: each root $v_0$ gives the straight line $y=v_0x$ through the origin. Write $v=y/x$ rather than $y=vx$ when the equation is presented in differential form and the combination $x\dd y-y\dd x$ is already visible.

\begin{worked}
$(x^{2}+y^{2})\dd x-2xy\dd y=0$, that is $y'=\dfrac{x^{2}+y^{2}}{2xy}$, with $y(1)=0$.

Both coefficients are degree $2$, so the quotient is degree $0$. With $y=vx$,
\[v+x\dv{v}{x}=\frac{1+v^{2}}{2v},\qquad x\dv{v}{x}=\frac{1-v^{2}}{2v}.\]
Separating, $\dfrac{2v\dd v}{1-v^{2}}=\dfrac{\dd x}{x}$, so $-\ln\abs{1-v^{2}}=\ln\abs x+c$ and $1-v^{2}=\dfrac{C}{x}$. Multiplying by $x^{2}$,
\[x^{2}-y^{2}=Cx.\]
With $y(1)=0$, $C=1$ and $y^{2}=x^{2}-x$; then $y(3)=\sqrt6=2.4494897$, confirmed numerically. The lost solutions are the roots of $F(v)-v=\dfrac{1-v^{2}}{2v}$, namely $v=\pm1$, giving the lines $y=\pm x$; and indeed $y=\pm x$ satisfies the original equation, as the case $C=0$.

A near miss worth studying: $x\dv{y}{x}-y=x^{2}+y^{2}$ is \emph{not} degree-zero homogeneous, since $y'=\dfrac{y}{x}+x+\dfrac{y^{2}}{x}$ mixes degrees. Yet $y=vx$ still works, because the substitution is really an assault on the combination $x\dd y-y\dd x$ rather than on homogeneity. Substituting gives $x^{2}v'=x^{2}(1+v^{2})$, so $v'=1+v^{2}$ exactly, $v=\tan(x+C)$, and $y=x\tan(x+C)$. With $y(\tfrac{\pi}{2})=\tfrac{\pi}{2}$ we need $C=-\tfrac{\pi}{4}$, so $y=x\tan\!\left(x-\tfrac{\pi}{4}\right)$, whence $y(\tfrac{\pi}{4})=0$ and $y(1.2)=0.5281360$, matching the numerics to seven figures. Equivalently, divide the original by $x^{2}$ to expose $\dd(y/x)=\bigl(1+(y/x)^{2}\bigr)\dd x$ and integrate at once. The moral: try $v=y/x$ whenever $x\dd y-y\dd x$ appears, whether or not the degree test passes.
\end{worked}

\begin{trap}
Writing $y'=v'$ instead of $y'=v+xv'$. This is the single most common error in the whole subject; it produces an equation that separates just as nicely and is simply wrong. The check costs nothing: differentiate $y=vx$ as a product. The second trap is discarding the roots of $F(v)-v$, which are the straight-line solutions through the origin and are usually the ones the geometry of the problem is about.
\end{trap}

\begin{seealsob}The homogeneous-of-degree-zero test; recognising a differential as $\dd$ of something; isobaric equations, the weighted generalisation.\end{seealsob}

It is worth pausing on the two ways of writing the same substitution, because they are convenient in different places. Writing $y=vx$ and differentiating gives $y'=v+xv'$, which is what you want when the equation is presented as $y'=F(y/x)$: substitute, cancel, separate. Writing $u=y/x$ instead and differentiating gives $\dd u=\dfrac{x\dd y-y\dd x}{x^{2}}$, which is what you want when the equation is presented in differential form and the combination $x\dd y-y\dd x$ is already sitting there. The two are the same substitution and produce the same separated equation; choosing the wrong one costs an extra line of algebra, choosing the right one sometimes finishes the problem outright, as the near-miss example above shows.

The second substitution handles the complementary situation: the right side does not depend on the ratio $y/x$ but on a fixed linear combination of $x$ and $y$. This family is easy to recognise because the combination is written out in the equation, usually inside a function.

\tech[Substitution u equals ax plus by plus c]{Substitution $u=ax+by+c$}{medium}
\cue The right side depends on $x$ and $y$ only through one linear combination: $y'=f(ax+by+c)$. Look inside every function in the equation; if $x$ and $y$ occur only in the pattern $ax+by+c$, and nowhere else, the substitution applies.
\method Set $u=ax+by+c$. Then $\dd u=a\dd x+b\dd y$, so
\[\dv{u}{x}=a+b\dv{y}{x}=a+bf(u),\]
which is autonomous in $u$ and therefore always separable. Solve for $u$, then recover $y=(u-ax-c)/b$. Translate the initial condition into $u$ at the same moment as the substitution.

\begin{worked}
$\dv{y}{x}=\cos(x+y)$ with $y(0)=0$. Here $u=x+y$, so $u'=1+\cos u=2\cos^{2}\tfrac u2$ by the half-angle identity, and
\[\tfrac12\sec^{2}\tfrac u2\dd u=\dd x,\qquad \tan\tfrac u2=x+C.\]
The condition $y(0)=0$ gives $u(0)=0$ and $C=0$, so $u=2\arctan x$ and
\[y=2\arctan x-x.\]
Check: $y'=\dfrac{2}{1+x^{2}}-1=\dfrac{1-x^{2}}{1+x^{2}}$, while $\cos(x+y)=\cos(2\arctan x)=\dfrac{1-x^{2}}{1+x^{2}}$. They agree, and $y(2)=2\arctan2-2=0.2142974$ matches the numerics. The lost solutions are the roots of $1+\cos u=0$, that is $u=(2k+1)\pi$, giving the straight lines $y=(2k+1)\pi-x$; these satisfy the original equation since $\cos(x+y)=-1=y'$ on each.

A second instance with a radical: $y'=\sqrt{x+y+1}$ with $y(0)=0$. With $u=x+y+1$, $u'=1+\sqrt u$. Substituting $s=\sqrt u$ makes the integral elementary, $\displaystyle\int\frac{2s\dd s}{1+s}=2s-2\ln(1+s)$, so
\[2\sqrt{u}-2\ln\bigl(1+\sqrt u\bigr)=x+C,\qquad C=2-2\ln2=0.6137056\]
from $u(0)=1$. This is implicit and should be left so; at $x=1$ the numerics give $y=1.4463061$, and substituting $u=x+y+1=3.4463061$ into the left side returns $x+C=1.6137056$ exactly as required.
\end{worked}

A third instance, where the combination is a ratio rather than an argument: $\dv{y}{x}=\dfrac{x+y-1}{x+y+1}$ with $y(0)=1$. Both numerator and denominator involve $x$ and $y$ only through $u=x+y$, so the substitution applies with $a=b=1$:
\[\dv{u}{x}=1+\frac{u-1}{u+1}=\frac{2u}{u+1},\qquad \frac{u+1}{2u}\dd u=\dd x,\]
so $u+\ln\abs u=2x+C$. From $u(0)=1$, $C=1$, and the solution is the implicit relation
\[x+y+\ln(x+y)=2x+1,\qquad\text{that is}\qquad y-x+\ln(x+y)=1.\]
At $x=1$ numerical integration gives $y=1.2079400$, and $u+\ln u=3.0000000$ as required. The lost solution is the root of $2u=0$, namely $u\equiv0$, that is the line $y=-x$; and indeed on that line the original right-hand side is $\dfrac{-1}{1}=-1$, which is the slope of $y=-x$.

\begin{trap}
Forgetting the $+a$ that comes from differentiating the $ax$ part, so that $u'=bf(u)$ instead of $u'=a+bf(u)$. The resulting equation is still separable and still solvable, so nothing warns you; the answer is simply wrong. The second trap is applying the substitution when $x$ and $y$ also appear outside the combination, for instance in $y'=x\cos(x+y)$, where $u'=1+x\cos u$ is not separable and the substitution has bought nothing.
\end{trap}

\begin{seealsob}Differentials as substitution bookkeeping; the substitution $y=vx$; autonomous recognition.\end{seealsob}

The two substitutions divide the territory neatly, and knowing which one applies is a matter of asking how $x$ and $y$ enter. If they enter only through the ratio $y/x$, use $y=vx$. If they enter only through a linear combination $ax+by+c$, use $u=ax+by+c$. If they enter through both, or through neither, neither substitution helps and you are back at the audit. The degenerate case $b=0$ is worth noting: then the equation is $y'=f(ax+c)$, a bare quadrature, and no substitution is needed at all.

The three substitutions covered so far exhaust the standard repertoire for hidden separability, and it is worth having them side by side with the question each answers.

\begin{center}
\footnotesize
\setlength{\tabcolsep}{5pt}
\renewcommand{\arraystretch}{1.42}
\begin{tabular}{@{}llll@{}}
\toprule
\mh{How $x$ and $y$ enter} & \mh{Substitute} & \mh{Becomes} & \mh{Lost solutions} \\
\midrule
Already a product $g(x)h(y)$ & nothing & $\dfrac{\dd y}{h}=g\dd x$ & roots of $h$ \\
Only through $y/x$ & $y=vx$ & $x v'=F(v)-v$ & lines $y=v_0x$ \\
Only through $ax+by+c$ & $u=ax+by+c$ & $u'=a+bf(u)$ & lines $ax+by+c=u_0$ \\
As $x\dd y-y\dd x$ over $x^{2}$ & $u=y/x$ & $\dd u=(\cdots)\dd x$ & $x=0$ \\
Only through $y x^{-m}$ & $y=vx^{m}$ & $x v'=G(v)-mv$ & curves $y=v_0x^{m}$ \\
\bottomrule
\end{tabular}
\end{center}

The last row is the weighted generalisation, and it is included to make the pattern visible: each substitution is chosen so that the combination through which $x$ and $y$ enter becomes a single new variable, after which the equation can only depend on that variable and on $x$, and separation follows. There is nothing to memorise beyond that principle. If you can name the combination, you can name the substitution.

\section{Trigonometric disguises}

The last family of disguises is the one where separation succeeds but the resulting integral does not look elementary. The separated equation is genuinely separable; what is hidden is the antiderivative, and the substitutions that expose it are the standard rationalising ones from integration technique, imported wholesale. The extra difficulty, and the reason this family is rated hard, is that the resulting antiderivatives are multivalued and the initial condition has to be tracked through the branches.

\tech{Trigonometric substitutions that separate}{hard}
\cue A separated side containing $1+\sin^{2}y$, $\sqrt{1-y^{2}}$, $1+y^{2}$, $a+b\cos y$, or any expression that would call for a trigonometric or Weierstrass substitution if it were an ordinary integral.
\method Separate first, then treat each side as an integration problem in its own right. The standard choices: $t=\tan y$ for even powers of $\sin$ and $\cos$; $y=\sin\theta$ for $\sqrt{1-y^{2}}$; $t=\tan\tfrac y2$ for a general rational function of $\sin y$ and $\cos y$. The one to have by heart, because it appears constantly, is
\[\int\frac{\dd y}{1+\sin^{2}y}=\frac{1}{\sqrt2}\arctan\bigl(\sqrt2\tan y\bigr)+\text{const},\]
obtained by dividing numerator and denominator by $\cos^{2}y$ and setting $t=\tan y$. Translate the initial condition into the new variable at the same time as the integrand, so that no branch ambiguity survives.

\begin{worked}
$\dv{y}{x}=\bigl(1+\sin^{2}y\bigr)\cos x$ with $y(0)=0$.

Separate: $\dfrac{\dd y}{1+\sin^{2}y}=\cos x\dd x$. The left side integrates by $t=\tan y$ as above, so
\[\frac{1}{\sqrt2}\arctan\bigl(\sqrt2\tan y\bigr)=\sin x,\]
the constant being zero because $y(0)=0$ sends both sides to zero. Solving,
\[y=\arctan\!\left(\frac{\tan\bigl(\sqrt2\sin x\bigr)}{\sqrt2}\right).\]
This formula is valid for all real $x$, and the reason is worth seeing: the argument $\sqrt2\sin x$ never leaves $[-\sqrt2,\sqrt2]$ and $\sqrt2=1.41421<\tfrac{\pi}{2}=1.57080$, so $\tan$ is never evaluated at a pole. Numerically, $y(3)=0.1420605$ and $y(5)=-1.2717020$ from both the formula and a Runge-Kutta integration.

Change the forcing to $2\cos x$ and the picture changes completely: the argument becomes $2\sqrt2\sin x$, which passes $\tfrac{\pi}{2}$, and the antiderivative $\tfrac{1}{\sqrt2}\arctan(\sqrt2\tan y)$ must be replaced by its continuous version
\[F(y)=\frac{1}{\sqrt2}\left[\arctan\bigl(\sqrt2\tan y\bigr)+\pi\floorb{\frac{y}{\pi}+\frac12}\right],\]
which adds $\tfrac{\pi}{\sqrt2}=2.2214415$ at each odd multiple of $\tfrac{\pi}{2}$. Checked against numerical quadrature, $F(5)-F(0)=3.4779665$, whereas the naive single-branch expression reports $-0.9649$, a value that is not even positive although the integrand is.
\end{worked}

The other trigonometric route worth having is the one for a rational function of $\sin$ and $\cos$. If separation leaves $\displaystyle\int\frac{\dd y}{a+b\cos y}$, the Weierstrass substitution $t=\tan\tfrac y2$ turns it into a rational integral in $t$, since $\cos y=\dfrac{1-t^{2}}{1+t^{2}}$ and $\dd y=\dfrac{2\dd t}{1+t^{2}}$. The result splits on the sign of $a^{2}-b^{2}$: an arctangent when $\abs a>\abs b$, a logarithm when $\abs a<\abs b$, and a rational function when $\abs a=\abs b$. The last case is the one to watch, because it is where the constant solutions live: if $a+b\cos y$ vanishes for some $y_0$, then $y\equiv y_0$ is a constant solution of the differential equation and the integral diverges there, which is the analytic shadow of the same fact.

\begin{trap}
Branch bookkeeping. The function $\arctan(\sqrt2\tan y)$ jumps by $\pi$ at every odd multiple of $\tfrac{\pi}{2}$, so it is \emph{not} an antiderivative of $\dfrac{1}{1+\sin^{2}y}$ across those points, even though its derivative is correct between them. Any initial condition, or any solution trajectory, that crosses $y=\tfrac{\pi}{2}$ requires the corrected $F$ above. The symptom of the error is a solution that is discontinuous where the differential equation is perfectly smooth, which should never be believed.
\end{trap}

\begin{seealsob}The separable machine with an initial condition; substitution $u=ax+by+c$; the Weierstrass substitution.\end{seealsob}

The simplest member of this family is worth a separate look, because it is the one where the branch question is easiest to see and hardest to remember. When both sides separate into arctangents, the solution is an addition formula in disguise, and the addition formula has a pole.

\begin{worked}[Worked example: arctangent on both sides]
$\dv{y}{x}=\dfrac{1+y^{2}}{1+x^{2}}$ with $y(0)=1$.

Separating gives $\arctan y=\arctan x+C$, and $y(0)=1$ forces $C=\tfrac{\pi}{4}$. Taking the tangent of both sides and using the addition formula,
\[y=\tan\!\left(\arctan x+\frac{\pi}{4}\right)=\frac{x+1}{1-x}.\]
Check: $y(0.5)=3$ and $y(0.9)=19$, both confirmed numerically to eleven figures. The solution exists only on $(-\infty,1)$, blowing up at $x=1$, and the differential equation gives no hint of this: its right-hand side is bounded by $1+y^{2}$ and perfectly smooth everywhere. What produced the pole is the arctangent's range. As $x$ increases from $0$, $\arctan x+\tfrac{\pi}{4}$ increases toward $\tfrac{\pi}{2}+\tfrac{\pi}{4}$ and crosses $\tfrac{\pi}{2}$ at $x=1$, at which point $y$ has escaped to infinity. Every separation that produces an arctangent on the $y$ side carries this hazard, and the check is the same each time: find where the accumulated right-hand side reaches $\tfrac{\pi}{2}$.
\end{worked}

\begin{keynote}[Read the closed form for its domain]
A separable equation is smooth wherever $g$ and $h$ are, and its solutions can still cease to exist at a finite $x$. The differential equation never announces this; the closed form always does, at the poles of a tangent, the zeros of a denominator, or the branch points of a logarithm. Three instances from this chapter: $y=\tfrac23\tan\bigl(3(1-\eu^{-(2x+1)})\bigr)$ blows up at $x^{*}=-0.1292526$; $y=1/(4x-x^{2})$ from the Bernoulli family exists only on $(0,4)$; and $y=x\tan(x-\tfrac{\pi}{4})$ exists only on $\bigl(-\tfrac{\pi}{4},\tfrac{3\pi}{4}\bigr)$. In every case the endpoints are invisible in the equation and explicit in the answer. Make reading the domain off the closed form the last step of every separation, and state it.
\end{keynote}

\section{Notation disguises}

A final category, cheap to handle once seen and paralysing when not. Some problems hide separability not in the algebra but in the typography, presenting the derivative in a form that has to be decoded before any classification is possible. The response is mechanical: translate to $\dv{y}{x}$ first, then run the audit as usual.

\tech[Notation disguises for dy dx]{Notation disguises for $\dv{y}{x}$}{easy}
\cue The derivative is buried in unusual notation: a $\dd x$-th root, a stacked exponential of differentials, a written-out limit of a difference quotient, or a phrase such as ``relative growth rate'' or ``proportional rate of change''.
\method Translate mechanically before doing anything else. The standard decodings: $\sqrt[\dd x]{\eu^{\dd y}}=\bigl(\eu^{\dd y}\bigr)^{1/\dd x}=\eu^{\dd y/\dd x}$; the relative growth rate $\dfrac{1}{M}\dv{M}{\theta}=\dvop{\theta}\ln M$; ``the rate of change of $y$ is proportional to $y$'' is $y'=ky$; and $\lim_{h\to0}\dfrac{f(x+h)-f(x)}{h}$ is $f'(x)$ however elaborately it is dressed. Then classify normally.

\begin{worked}
Solve $x\eu^{\eu^{x}}=\sqrt[\dd x]{\eu^{\dd y}}$ with $y(1)=\eu$, and find $y(\eu)$.

Decode the right side: it is $\eu^{\dd y/\dd x}=\eu^{y'}$. Taking logarithms of both sides,
\[y'=\ln\bigl(x\eu^{\eu^{x}}\bigr)=\ln x+\eu^{x}.\]
This is a bare quadrature, the most degenerate separable case. Integrating, $y=x\ln x-x+\eu^{x}+C$, and $y(1)=\eu$ forces $-1+\eu+C=\eu$, so $C=1$. Hence
\[y(\eu)=\eu\cdot1-\eu+\eu^{\eu}+1=\eu^{\eu}+1=16.1542622,\]
confirmed by numerical integration to nine figures.

A second instance: a population satisfies $\dfrac{1}{M}\dv{M}{\theta}=\cos\theta$ with $M(0)=2$. The left side is $\dvop{\theta}\ln M$, so $\ln M=\sin\theta+\ln2$ and $M=2\eu^{\sin\theta}$: a bounded oscillation between $2\eu^{-1}$ and $2\eu$, not the runaway growth that the word ``growth rate'' suggests to a careless reader.
\end{worked}

\begin{trap}
Treating $\dd y$ and $\dd x$ as separate objects that can be manipulated outside a quotient. The disguise works precisely because it invites that: writing $\sqrt[\dd x]{\eu^{\dd y}}$ tempts you to take a $\dd x$-th root of something, when in fact the only meaningful object present is the ratio $\dd y/\dd x$. Decode to a ratio first and the temptation disappears.
\end{trap}

\begin{seealsob}Differentials as substitution bookkeeping; hidden separability by factoring; the classification audit.\end{seealsob}

One further disguise belongs here even though it is verbal rather than typographic. Applied problems state the differential equation in words, and the words have standard translations that are worth running mechanically rather than interpreting. ``The rate of change of $y$ is proportional to $y$'' is $y'=ky$. ``Proportional to the difference between $y$ and the ambient value $A$'' is $y'=k(A-y)$, Newton's law of cooling, which is autonomous with the single equilibrium $y=A$. ``Proportional to the product of the number infected and the number not yet infected'' is $y'=ky(N-y)$, logistic. ``Inversely proportional to $y$'' is $y'=k/y$, which separates to $y^{2}=2kx+C$. In each case the constant $k$ carries a sign that the wording fixes and the algebra does not, and getting that sign wrong is the commonest error in the whole applied family: cooling has $k>0$ with the difference written $A-y$, and writing $y-A$ instead flips the stability of the equilibrium without changing anything else in the algebra.

\section{What separation cannot do}

Two limits are worth stating plainly, because both are places where a correct separation produces a wrong-looking answer and the temptation is to go back and re-derive.

The first is that separation frequently produces an implicit relation which cannot be solved for $y$, and that this is a finished answer rather than a stalled one. The radical example above, $2\sqrt u-2\ln(1+\sqrt u)=x+C$ with $u=x+y+1$, is of this kind: no rearrangement extracts $y$, and none is needed, because every question one can reasonably ask (the value of $y$ at a given $x$, the location of a vertical tangent, the behaviour as $x\to\infty$) is answerable from the implicit form by inspection or by one numerical root-find. Solving for $y$ is a convenience, not a requirement, and forcing it usually costs a case split on signs that the implicit form did not need.

The second is that the antiderivatives which appear after separation need not be elementary. The equation $y'=\eu^{-x^{2}}(1+y^{2})$ separates perfectly into $\dfrac{\dd y}{1+y^{2}}=\eu^{-x^{2}}\dd x$, and the right-hand side is $\tfrac{\sqrt\pi}{2}\erf(x)$, which is not elementary. The correct response is to write $\arctan y=\tfrac{\sqrt\pi}{2}\erf(x)+C$ and stop. Separation has done its entire job: it has reduced a differential equation to two independent integrations. Whether those integrations have elementary answers is a separate question belonging to a different chapter, and a non-elementary antiderivative is not evidence that the separation was wrong.

Finally, a word on what the constant-solution discipline buys beyond correctness. Because solutions of a Lipschitz equation cannot cross, every constant solution is a barrier, and the barriers organise the whole solution family: they partition the strip in which the solutions live, and each non-constant solution is trapped between two consecutive barriers for its entire existence. That is why recording the roots of $h$ first pays twice, once as a completeness check on the answer and once as a free qualitative description of every solution you did not compute.

\section{Recognition matrix}
\begin{recog}{@{}p{0.31\textwidth}p{0.35\textwidth}p{0.28\textwidth}@{}}
\rowcolor{mist}\mh{If you see} & \mh{Reach for} & \mh{If that fails} \\
\midrule
\endhead
$y'=g(x)h(y)$ with data $y(x_0)=y_0$ & Definite form $\int_{y_0}^{y}\dd s/h=\int_{x_0}^{x}g$ & Indefinite form, then fit $C$ \\
Any division by $h(y)$ & Record the roots of $h$ first & Check them against the final formula \\
$\ln\abs y$ after integrating & $A=\pm\eu^{C}$ is any nonzero real; $A=0$ is the equilibrium & State the interval of validity \\
A four-term sum on the right & Group in pairs and factor & Test linearity in $y$ \\
$\eu^{x+y}$, $\eu^{x-y}$, $a^{x+y}$ & Split the exponential into a product & Substitute $u=x+y$ \\
A radicand like $x+2\sqrt{x-1}$ & Complete the square inside; watch $\sqrt{u^{2}}=\abs u$ & State the domain explicitly \\
Trigonometric coefficient of two or more terms & Apply a Pythagorean identity first & Weierstrass $t=\tan(x/2)$ \\
$f(tx,ty)=f(x,y)$ & $y=vx$, then $x v'=F(v)-v$ & Check $y'=v+xv'$, not $v'$ \\
$x\dd y-y\dd x$ visible & $v=y/x$ directly; divide by $x^{2}$ & Divide by $x^{2}+y^{2}$ for $\arctan$ \\
Roots of $F(v)-v$ & Straight-line solutions $y=v_0x$; report them & They are the $C=0$ case \\
$y'=f(ax+by+c)$ & $u=ax+by+c$; $u'=a+bf(u)$ & Do not drop the $+a$ \\
$x$ or $y$ also outside the combination & The substitution fails; audit again & Try $y=vx$ or exactness \\
$1+\sin^{2}y$ downstairs & $t=\tan y$; get $\tfrac{1}{\sqrt2}\arctan(\sqrt2\tan y)$ & Add $\pi\floorb{y/\pi+\tfrac12}$ for continuity \\
$\sqrt{1-y^{2}}$ on the $y$ side & $y=\sin\theta$; watch the range of $\arcsin$ & Separate and leave implicit \\
Rational in $\sin y$ and $\cos y$ & $t=\tan(y/2)$ & Check for a hidden constant solution \\
A $\dd x$-th root or stacked differentials & Decode to $\dd y/\dd x$ first, then classify & Take logarithms of both sides \\
``Relative growth rate'' & $\dfrac1M\dv{M}{\theta}=\dvop{\theta}\ln M$ & Separate in $\ln M$ \\
A closed form with $\tan$ or a denominator & Locate its poles: the solution may end there & Compare against a numerical run \\
\end{recog}
